\documentclass[sigconf]{acmart}

\title{From Resume Drafting to Interview Rehearsal: Understanding How University Students Use LLMs for Career Preparation}

\author{Anonymous Authors}
\affiliation{%
  \institution{Anonymous Institution}
  \city{Anonymous City}
  \country{Anonymous Country}
}
\email{anonymous@example.com}

\begin{document}

\begin{abstract}
Large language models have quickly entered student career-preparation workflows, where they are used to draft resumes, rehearse interview answers, compare options, and reduce uncertainty during job search. Yet existing literature does not clearly explain how these systems are embedded across the broader ecology of career preparation, nor how students interpret their influence on confidence, perceived employability, authenticity, and skill development. This manuscript draft frames the project as a CHI study of LLM-mediated career preparation rather than a claim about objective employment outcomes. To validate the local-first analysis and writing pipeline before live deployment, we generated a synthetic dataset matching the planned questionnaire schema and analyzed it as a dry run. The synthetic results suggest a plausible pattern in which LLM use concentrates around resume drafting, portfolio polishing, and interview rehearsal, while concerns about authenticity, deskilling, and safeguard needs remain prominent. The resulting contribution is a reproducible end-to-end workflow for studying student-facing AI career-support systems and a draft empirical narrative ready to be replaced with real data once collection begins.
\end{abstract}

\maketitle

\section{Introduction}
Public discussion about generative AI often collapses a complicated set of educational and labor-market changes into a simple claim that large language models will either help or harm student employment. That framing is too coarse for human-computer interaction. What matters for HCI is not whether AI changes macro employment in the abstract, but how students actually use LLMs while preparing resumes, researching roles, rehearsing interviews, comparing pathways, and deciding how much of their own voice to preserve. Prior work has already shown that generative AI introduces new usability and metacognitive demands around prompting, evaluating outputs, and calibrating trust \cite{tankelevitch2024metacognitive}. At the same time, higher-education research suggests that students perceive both substantial value and substantial risk in these tools, including concerns about accuracy, privacy, authenticity, and long-term development \cite{chan2023studentsvoices,jochim2024teachingtesting}.

The career-preparation setting raises these tensions sharply. A polished answer may increase confidence while simultaneously reducing reflection. A rewritten resume bullet may sound stronger while making the applicant less certain that the application still represents them. Existing work on student use of generative AI has not yet fully addressed this setting. Some studies focus on specific populations such as students with disabilities \cite{atcheson2025disabilitygenai}; others explore future educational applications through co-design \cite{prasad2025multimodal}; still others present narrow technical systems such as AI-supported interview simulators \cite{lee2024vrinterview}. The closest topical study reports a positive association between ChatGPT use and students' employment confidence \cite{xiao2025employmentconfidence}, but it does not unpack the concrete interaction practices through which confidence is built, misplaced, or negotiated.

This project therefore asks how university students use LLMs across career-preparation tasks, how they perceive those interactions as shaping confidence and employability, and what design implications follow for responsible student-facing AI career support. The aim is to produce a grounded HCI account of LLM-mediated career preparation while explicitly avoiding unsupported causal claims about objective labor-market outcomes.

\section{Related Work}
Our starting point is the growing HCI literature on the demands of generative AI use. Tankelevitch et al.\ argue that prompting, checking, and deciding whether to rely on outputs are fundamentally metacognitive activities, which makes generative AI both useful and cognitively demanding in real work \cite{tankelevitch2024metacognitive}. Schemmer et al.\ similarly show that appropriate reliance on AI advice depends on how explanations shape users' calibration and trust \cite{schemmer2023appropriatereliance}. That framing is directly relevant to career preparation, where students must decide not only whether an output is fluent, but whether it is accurate, appropriate, and representative of their abilities.

We also build on research about students' experiences with generative AI in higher education. Chan and Hu report that students are broadly optimistic about generative AI while also expressing concerns about accuracy, ethics, privacy, and career prospects \cite{chan2023studentsvoices}. Jochim and Lenz-Kesekamp similarly describe adaptation pressures on both students and teachers in the era of text-generative AI \cite{jochim2024teachingtesting}. Atcheson et al.\ add a more specific and important lens by showing that students with disabilities use generative AI not only for coursework, but also for self-advocacy and support, while simultaneously navigating unclear institutional norms \cite{atcheson2025disabilitygenai}. Prasad et al.\ show through co-design workshops that students imagine multimodal generative AI tools as personalized educational support systems \cite{prasad2025multimodal}. These studies establish the importance of student perspectives, but they are not centered on career preparation as a distinct and high-stakes interaction context.

Finally, adjacent career-oriented systems suggest a fragmented design space. Lee et al.\ present a VR and chatbot-based interview simulator for graduates in Hong Kong \cite{lee2024vrinterview}, and Waikar et al.\ describe a generative-AI career-guidance system aimed at student support \cite{waikar2024careerguidance}. Xiao and Zheng report that ChatGPT use is associated with higher self-reported employment confidence \cite{xiao2025employmentconfidence}. Meanwhile, Chakrabarty et al.\ warn that LLMs can create a false promise of creativity and authorship, a concern that maps directly onto resume writing and personal statement preparation \cite{chakrabarty2024artorartifice}. Together, these works indicate that AI is already entering student career support, but they do not provide a broad account of how students stitch LLMs into everyday career-preparation workflows or where design interventions should preserve voice, promote verification, and reduce overreliance. That is the gap this project addresses.

\section{Method}
The study is designed as a survey-first mixed-method investigation. The main data source is an online questionnaire administered to university students and recent graduates who are actively preparing for internships, full-time jobs, graduate school, or adjacent early-career transitions. The survey captures whether and how respondents use LLMs for concrete tasks such as resume drafting, cover-letter writing, interview rehearsal, career exploration, networking-message drafting, portfolio polishing, and option comparison. It also measures perceived confidence support, employability support, verification behavior, authenticity concerns, and perceived deskilling risk. Conceptually, the instrument combines CHI-motivated interaction constructs with adjacent career-development measures related to perceived employability and career self-efficacy \cite{bennett2022employability,jin2009careerdecision}.

The instrument is intentionally task-specific because broad opinion questions about AI tend to conflate hype, anxiety, and actual practice. Respondents are therefore asked about concrete use frequencies, typical interaction patterns, scenario judgments, and open-ended examples of both useful and risky encounters. The project retains both users and non-users so that adoption barriers and contrast cases can be studied alongside active reliance. Data collection is planned through a survey platform such as Qualtrics or Wenjuanxing, while the project repository remains the local source of truth for the instrument, export schema, and analysis plan.

At the end of the survey, respondents may optionally volunteer for a short follow-up interview. These interviews are intended to explain mechanisms that are difficult to capture through scales alone, such as how students decide whether an LLM-generated answer still sounds like them, when AI reduces uncertainty versus when it replaces reflection, and what safeguards they want from future AI career-support systems. Contact information for interviews is stored separately from main survey responses to preserve privacy.

Because real data collection had not yet begun in this local-first run, we generated a synthetic survey export that matched the planned schema and approximate study distributions. The synthetic dataset contained 336 raw rows and was analyzed using the same exclusion logic, descriptive statistics, and regression pipeline that will later be applied to real exports. This choice does not create empirical evidence; instead, it allows the project to validate file formats, figure generation, manuscript structure, and claim boundaries before participant recruitment.

\section{Results}
All findings in this section are derived from a synthetic dataset and are reported only to exercise the analysis and writing pipeline. After applying the predefined exclusions, the cleaned synthetic sample contained 319 cases, of which 78.7\% were marked as LLM users. The retained sample was composed primarily of undergraduates (63.3\%) and master's students (22.3\%), with the largest career-preparation contexts being full-time employment (34.8\%) and internships (32.3\%).

Task-level synthetic usage followed an interpretable pattern rather than a uniform one. Among LLM users, the most common weekly-or-more activities were resume drafting (59.4\%), portfolio polishing (53.0\%), and interview rehearsal (51.8\%). Cover-letter writing, information search, and skill-gap planning formed a middle tier of use, while networking and offer comparison were less frequent. This pattern fits the broader claim that students may be especially likely to use LLMs in high-friction moments where language production, self-presentation, and iterative preparation are central.

The synthetic group contrasts suggest a graded relationship between LLM engagement and perceived support. Heavy users reported higher confidence-support scores (M = 3.116) and higher perceived employability-support scores (M = 3.807) than light users (M = 2.662 and 3.057) and non-users (M = 2.319 and 2.191). Verification scores remained relatively similar across these groups, whereas authenticity tension increased with heavier use. In the synthetic regression predicting perceived employability support, use intensity was positively associated with the outcome ($\beta = 0.4823$, $p < .001$), verification was also positively associated with it ($\beta = 0.1490$, $p = .029$), and authenticity tension was negatively associated with it ($\beta = -0.1821$, $p = .026$). In the parallel confidence-support model, use intensity remained a positive predictor ($\beta = 0.3110$, $p < .001$), while authenticity tension was effectively null in this dry run.

The synthetic open-text analysis converged on a small number of recurring themes. Desired safeguards appeared in the full sample by construction of the generated responses, while authenticity and voice concerns (73.0\%) and deskilling or dependency concerns (70.2\%) also appeared frequently. Verification and trust (51.1\%) and interview rehearsal (55.5\%) were similarly common. Taken together, these synthetic results produce a coherent placeholder story in which students use LLMs most for preparing and polishing self-presentation tasks, experience meaningful perceived support, and still want mechanisms that preserve authorship and encourage checking.

\section{Discussion}
Even though the present findings are synthetic, the dry run is already useful for sharpening the paper's eventual argument. Career preparation is not a single interaction but a chain of situated tasks in which students repeatedly negotiate speed, confidence, authorship, and trust. The placeholder results reinforce the likelihood that LLMs will be especially attractive in self-presentation tasks such as resume drafting, portfolio polishing, and interview rehearsal, where immediate fluency gains are valuable. At the same time, the synthetic theme structure suggests that adoption is unlikely to be a simple story of benefit. Students may accept help precisely where they also most fear losing authorship or overfitting themselves to generic AI output.

This points to four concrete design implications for future student-facing systems. First, systems should preserve authorship by helping students compare their own draft with model-generated alternatives rather than replacing their text outright. Second, systems should support verification by distinguishing language polishing from substantive career advice and by prompting external checks for high-stakes recommendations, consistent with prior work on calibrated reliance \cite{schemmer2023appropriatereliance}. Third, systems should help students reflect on skill development by making visible what was changed and what still requires independent practice. Fourth, career centers and educational institutions may need clearer guidance for acceptable and unacceptable forms of LLM use during application and interview preparation. These themes connect everyday student practice to a more responsible design agenda for AI-supported career development.

\section{Limitations}
The strongest limitation is that the current analysis uses synthetic rather than empirical data. The numerical results are therefore not evidence about real students and must be treated only as placeholders for validating the pipeline, tables, figures, and prose structure. Once real data are collected, the study will still face important limitations. The design is centered on self-report and cannot establish that LLM use objectively improves employment outcomes. Early recruitment is likely to be convenience-based, which may overrepresent students who are already digitally engaged or especially opinionated about AI. If the first collection wave is China-first, the cultural and labor-market context will also shape what can be generalized beyond that setting. These limitations reinforce the decision to keep the paper's claims focused on practices, perceptions, and design implications.

\section{Ethics and Privacy}
The study collects minimal personal data and does not require participants to upload resumes, employer names, or offer materials. Survey responses are analyzed in aggregate, while interview contact information is separated from the main dataset. The project treats AI use in career preparation as a potentially sensitive topic because respondents may worry about judgment, cheating narratives, or perceived lack of competence. For that reason, the instrument uses neutral wording, allows participants to skip questions, and avoids framing LLM use as inherently good or bad. Any eventual publication will state clearly that the data speak to reported practice and perceived confidence, not to causal changes in objective employment outcomes.

\section{Conclusion}
This project contributes a grounded HCI account of how university students integrate LLMs into career-preparation workflows and how those practices shape perceived confidence, perceived employability support, and tensions around authenticity and skill development. By reframing the original question --- from ``do LLMs affect employment?'' to ``how do students use LLMs during career preparation, and what should responsible systems do?'' --- we position the work at a level where HCI methods and design implications are most actionable.

The local-first pipeline developed here demonstrates that a complete survey instrument, analysis plan, and draft manuscript can be validated end-to-end using synthetic data before a single participant is recruited. The four design implications --- authorship preservation, verification scaffolding, skill-development visibility, and institutional guidance --- emerge from the synthetic dry run as candidate design directions that real data can either confirm or complicate. The next step is a piloted deployment of the survey instrument, replacement of synthetic placeholders with live data, and refinement of the discussion and contribution claims based on actual student voices.

More broadly, this project illustrates the kind of local-first research infrastructure that CHI HCI work on student-facing AI systems can follow: maintain the instrument, schema, analysis scripts, and manuscript in a single reproducible project directory; use synthetic dry runs to validate the full pipeline before the study opens; and preserve a clear boundary between empirical claims, design implications, and causal speculation. We hope that future work on AI-supported career development can build directly on the open materials in this project.

\bibliographystyle{ACM-Reference-Format}
\bibliography{references}

\end{document}
